\section{What these bounds do, and do not, imply asymptotically}
\label{sec:complexity}

\subsection{A complete collapse-only trajectory is a polynomial subroutine}

Let $V(t,B)$ bound the cost of the maintained finite-manifold and cocycle
preparations on $t$ tetrahedra with $B$-bit edge differences. Initial
height rows are assumed gauged, for example by subtracting their first
entry. If arbitrarily large common offsets are supplied, reading and
normalising those integers is charged separately in their actual input
bit length. We keep the preparation cost explicit. Let $\M(B)$ stand for a safe bound on the relevant exact
integer arithmetic cost; additions and comparisons are cheaper than general
multiplication. A complete current census has $O(t)$ candidates, because
there are at most $6t$ local edge occurrences. Sorting for greedy packing
costs $O(t\log(t+2))$ comparisons. The minima index and source-specific
matching add polynomial work independent of expanded surface weight.

\begin{proposition}[Cost and termination of the supplied local policy]
\label{prop:collapse-trajectory}
Start with a coherent cocycle on a valid triangulation with at most two
interior vertices. At each round, form the complete legal $3$--$2$ census,
choose a greedy maximal disjoint family, select its optimal peeled-Euler
subbatch with the minimum-omission rule, and commit it if nonempty. Stop
otherwise. Every committed round is nondecreasing in peeled Euler
characteristic; there are at most $t_0$ rounds and at most $t_0$ individual
collapses, where $t_0$ is the initial tetrahedron count. The maximum encoded
edge-difference bit length does not increase.

With polynomial $V$ and exact arithmetic, this is a polynomial subroutine
in $(t_0,B_0)$. A deliberately loose explicit bound is
\begin{equation}\label{eq:trajectory-cost}
 O\!\left(t_0 V(t_0,B_0)+
 t_0^2\log(t_0+2)\,\M(B_0+\log(t_0+2))\right),
\end{equation}
assuming a nondecreasing upper bound $V$. The lifetime maintenance work of
the corner index alone is only
$O(t_0\log(t_0+2))$ record operations.
\end{proposition}
\begin{proof}
The empty subbatch is available and has gain zero, so an optimum cannot
have negative gain. Every committed move reduces the number of tetrahedra
by one. Thus at most $t_0$ moves can be committed, and each nonempty round
contains at least one. The commuting theorem preserves the vertex types,
so the small-cover selection remains applicable. Every edge after a
$3$--$2$ move already occurred on the bipyramid boundary; its cocycle value
is unchanged. This proves the bit statement.

Charge each round its complete census, packing, reward computation,
linear small-cover matching, final emission, and constant number of
preparations. Bounding every current size by $t_0$ and the number of rounds
by $t_0$ gives \eqref{eq:trajectory-cost}. There are initially $4t_0$ corner
records; each collapse deletes twelve and inserts eight, so the total
number of index edits is $O(t_0)$. The balanced-tree theorem supplies the
last assertion.
\end{proof}

This proposition gives an honest finite local algorithm, including its
stopping condition. It makes no claim that the terminal coherent vector
contains a disc or that a stalled triangulation cannot be the unknot.
Repeating an uncontrolled $2$--$3$ repair phase is outside its proof.
Likewise, recording every endpoint of every round can still cost quadratic
space in $t_0$; the per-batch encoding theorem does not automatically supply
a linear-size transcript for an entire trajectory.

\subsection{A density hypothesis would give logarithmically many rounds}

\begin{corollary}[Conditional dense-census progress]
\label{cor:dense-progress}
Suppose a collapse-only trajectory as above has at least $\alpha t$
eligible distinct central edges whenever its current size is $t\ge40/\alpha$,
where $0<\alpha\le6$ is fixed. Then it reaches size below $40/\alpha$ in
$O(\alpha^{-1}\log(t_0+2))$ committed rounds.
\end{corollary}
\begin{proof}
The packed optimum retains at least
\[
 \lceil\alpha t/10\rceil-2\ge\alpha t/20
 \qquad(t\ge40/\alpha).
\]
Thus the next size is at most $(1-\alpha/20)t$. Repeated multiplication
proves the bound, using $-\log(1-\alpha/20)\ge\alpha/20$.
\end{proof}

The density assumption is not established for all diagram exteriors or
after arbitrary moves. A legal census can be sparse or empty. The corollary
is useful because it specifies an exact geometric statement that would
turn the batch-retention theorem into a round bound, without conflating
raw score, peeled score, and the number of removed tetrahedra.

\subsection{The normal-search ledger}

The active-face theorem isolates the factor
$\mathcal A(T)2^{\rho_{\max}}$. It supplies no universal reduction of this
factor. In fact, the unrestricted source lower bound rules out the proposal
that $\rho_{\max}$ is automatically polylogarithmic in crossing number.
The anchor count also needs separate attention. In a one-vertex finite
triangulation it is $4t$; in a triangulation with many vertices it is the
product of their corner degrees. The fact that there are only two
\emph{interior} vertices in the pulling source does not make that anchor
product small: the same source has $4n'$ boundary vertices. Confusing the
two vertex counts would erase a genuine source of combinatorial cost.

Complete-support certification, support minimisation of one admissible
point, and proving a small active ambiguity rank are three different
tasks. The first now has a compact interface. The second already has
classical LP-based solutions \cite{BurtonOzlen}. The third fails at the
unrestricted roots proved above. A promising parameter must therefore be
defined after stronger certified restrictions, or the algorithm must be
proved to reach an endpoint before traversing the large rank.

\subsection{A sufficient quasipolynomial contract}

For clarity, consider a future complete architecture with at most $R(n)$
outer stages. At each stage suppose it covers all required geometric
choices by a tree of depth $h(n)$, branching at most $(ct(n)+c)^d$ ways per
node for fixed positive constants $c,d$, and each node performs at most $A(n)$ normal queries with active
rank at most $\rho(n)$. Suppose all state sizes are bounded by $t(n)$ and
all arithmetic inputs by $B(n)$ bits, with polynomial local producers and
consumers. A sufficient upper bound is
\begin{equation}\label{eq:qp-ledger}
 R(n) A(n) 2^{\rho(n)}
 (ct(n)+c)^{O(h(n))}\poly(t(n),B(n)).
\end{equation}
For example, if $t,B,R$ are polynomial in $n$,
$\log A(n)=O((\log n)^a)$, $\rho(n)=O((\log n)^b)$, and
$h(n)=O((\log n)^c)$ for fixed exponents, then
\eqref{eq:qp-ledger} is
\[
 \exp\!\left(O((\log n)^{\max\{1,a,b,c+1\}})\right).
\]
This is an elementary composition bound, not a theorem that those
hypotheses hold. It displays the obligations a genuine quasipolynomial
proof must discharge.

The present results improve the local polynomial factor and eliminate
$2^{|E|}$ subbatch enumeration for a fixed disjoint family. They do not
prove geometric witness coverage, the depth bound $h$, the restricted-rank
bound $\rho$, or the anchor bound $A$. The explicit root obstruction also
shows why replacing the last two obligations by an empirical observation
would not suffice. Classical compressed component algorithms avoid
dependence linear in expanded normal weight \cite{AHT}; they do not bound
the number of global choices by themselves.

